% !TEX root = ../main.tex
% ============================================================================
%                           PENGATURAN HALAMAN
% ============================================================================
% Margin sesuai panduan UDINUS:
% - Kiri: 4 cm (untuk jilid)
% - Kanan: 2.5 cm
% - Atas: 3 cm
% - Bawah: 3 cm
\usepackage[a4paper,left=4cm,right=2.5cm,top=3cm,bottom=3cm]{geometry}
\usepackage{tabularx}
\usepackage{xltabular}

% ============================================================================
%                           PENGATURAN FONT
% ============================================================================
% Font wajib: Times New Roman 12pt (sesuai panduan)
\usepackage{fontspec}
\setmainfont{Times New Roman}
\usepackage{unicode-math}
\setmathfont{Latin Modern Math}

% ============================================================================
%                           WARNA KHUSUS COVER
% ============================================================================
\usepackage{xcolor}
\definecolor{darkred}{RGB}{139,0,0}  % Warna merah gelap untuk cover
\definecolor{codegreen}{rgb}{0,0.6,0}
\definecolor{codegray}{rgb}{0.5,0.5,0.5}
\definecolor{codepurple}{rgb}{0.58,0,0.82}
\definecolor{codeblue}{rgb}{0.13,0.29,0.53}
\definecolor{codeorange}{rgb}{0.8,0.4,0}
\definecolor{backcolour}{rgb}{0.97,0.97,0.97}

% ============================================================================
%                           PENGATURAN SPASI
% ============================================================================
% Spasi 1.5 sesuai panduan penulisan tesis
\usepackage{setspace}
\onehalfspacing
\usepackage{indentfirst}  % Indent paragraf pertama setiap section
\usepackage{enumitem}
\setlist[enumerate]{leftmargin=*, itemsep=0.5em}
\setlist[itemize]{leftmargin=*, itemsep=0.5em}

% ============================================================================
%                           PENGATURAN BAHASA
% ============================================================================
\usepackage[english]{babel}
\babelprovide[main]{indonesian}


% ============================================================================
%                        HEADER & FOOTER (NOMOR HALAMAN)
% ============================================================================
\usepackage{fancyhdr}

% ============================================================================
%                        FORMAT JUDUL BAB & SUBBAB
% ============================================================================
\usepackage{titlesec}
\usepackage{etoolbox}

% ============================================================================
%                         KONFIGURASI DAFTAR ISI (TOC)
% ============================================================================
\usepackage{tocloft}

% ============================================================================
%                    GAMBAR, TABEL, DAN CAPTION
% ============================================================================
\usepackage{graphicx}
\usepackage{pdfpages}  % Untuk \includepdf pada lampiran
\usepackage[export]{adjustbox}  % max width=\textwidth untuk tabel/gambar overfull
\usepackage{caption}

% --- TikZ untuk diagram dan flowchart ---
\usepackage{tikz}
\usepackage{pgfplots}
\pgfplotsset{compat=1.18}
\usetikzlibrary{shapes.geometric, shapes.misc, arrows.meta, positioning, calc, decorations.pathreplacing, backgrounds, fit}

% Pastikan amsmath dimuat untuk kontrol persamaan
\usepackage{amsmath}

% ============================================================================
%                         PSEUDOCODE/ALGORITMA
% ============================================================================
\usepackage[ruled,vlined,linesnumbered,resetcount]{algorithm2e}
\SetKwInput{KwInput}{Input}
\SetKwInput{KwOutput}{Output}
\SetKw{KwAnd}{and}
\SetKw{KwOr}{or}
\SetAlgoCaptionSeparator{.}
\SetAlgoNlRelativeSize{-1}
\SetAlFnt{\small}
\SetKwComment{Comment}{$\triangleright$ }{}

% ============================================================================
%                         KODE SUMBER (SOURCE CODE)
% ============================================================================
\usepackage{listings}

% --- Style Default untuk Kode ---
\lstdefinestyle{defaultstyle}{
    backgroundcolor=\color{backcolour},
    commentstyle=\color{codegreen}\itshape,
    keywordstyle=\color{codeblue}\bfseries,
    numberstyle=\tiny\color{codegray},
    stringstyle=\color{codeorange},
    basicstyle=\ttfamily\small,
    breakatwhitespace=false,
    breaklines=true,
    captionpos=b,
    keepspaces=true,
    numbers=left,
    numbersep=8pt,
    showspaces=false,
    showstringspaces=false,
    showtabs=false,
    tabsize=2,
    frame=single,
    framerule=0.5pt,
    rulecolor=\color{codegray},
    xleftmargin=2em,
    framexleftmargin=1.5em
}

% --- Style khusus Python ---
\lstdefinestyle{pythonstyle}{
    style=defaultstyle,
    language=Python,
    morekeywords={self, True, False, None, as, with, async, await, yield, lambda},
    emph={def, class, return, import, from, if, elif, else, for, while, try, except, finally, raise, pass, break, continue, and, or, not, in, is},
    emphstyle=\color{codepurple}\bfseries
}

% --- Style khusus JavaScript (definisi manual karena tidak built-in) ---
\lstdefinelanguage{JavaScript}{
    keywords={break, case, catch, continue, debugger, default, delete, do, else, finally, for, function, if, in, instanceof, new, return, switch, this, throw, try, typeof, var, void, while, with, class, const, enum, export, extends, import, super, implements, interface, let, package, private, protected, public, static, yield, async, await},
    morecomment=[l]{//},
    morecomment=[s]{/*}{*/},
    morestring=[b]',
    morestring=[b]",
    morestring=[b]`,
    sensitive=true
}

\lstdefinestyle{javascriptstyle}{
    style=defaultstyle,
    language=JavaScript,
    keywordstyle=\color{codeblue}\bfseries,
    commentstyle=\color{codegreen}\itshape,
    stringstyle=\color{codeorange}
}

% --- Style khusus Java ---
\lstdefinestyle{javastyle}{
    style=defaultstyle,
    language=Java,
    morekeywords={String, System, out, println},
    emph={public, private, protected, static, void, class, interface, extends, implements, new, return, if, else, for, while, try, catch, finally, throw},
    emphstyle=\color{codepurple}\bfseries
}

% --- Style khusus Prompt LLM ---
\lstdefinelanguage{PromptLLM}{
    moredelim=[s][\color{codeblue}\bfseries]{[}{]},
    moredelim=[s][\color{codeorange}\bfseries]{\{}{\}},
    morekeywords={Output, JSON, Return, keys, with},
    keywordstyle=\color{codepurple}\bfseries,
    morecomment=[l]{-},
    commentstyle=\color{codegreen},
}

\lstdefinestyle{promptstyle}{
    style=defaultstyle,
    language=PromptLLM,
    basicstyle=\ttfamily\scriptsize,
}

% Set style default
\lstset{style=defaultstyle}

% Dukungan karakter spesial (Greek, dll)
\lstset{
  extendedchars=true,
  literate={α}{{$\alpha$}}1 {γ}{{$\gamma$}}1 {ε}{{$\epsilon$}}1
}

% Counter sharing: lstlisting dan algocf pakai nomor urut yang sama, keduanya masuk .loa
\makeatletter
\@addtoreset{lstlisting}{chapter}
\renewcommand{\ext@lstlisting}{loa}

\let\orig@refstepcounter\refstepcounter
\renewcommand{\refstepcounter}[1]{%
  \def\temp@arg{#1}%
  \def\temp@lst{lstlisting}%
  \ifx\temp@arg\temp@lst
    \orig@refstepcounter{algocf}%
    \setcounter{lstlisting}{\value{algocf}}%
  \else
    \orig@refstepcounter{#1}%
  \fi
}

\let\orig@stepcounter\stepcounter
\renewcommand{\stepcounter}[1]{%
  \def\temp@arg{#1}%
  \def\temp@lst{lstlisting}%
  \ifx\temp@arg\temp@lst
    \orig@stepcounter{algocf}%
    \setcounter{lstlisting}{\value{algocf}}%
  \else
    \orig@stepcounter{#1}%
  \fi
}

\AtBeginDocument{
  \@addtoreset{algocf}{chapter}
  \@addtoreset{lstlisting}{chapter}
  \renewcommand{\thelstlisting}{\thealgocf}
  \let\l@lstlisting\l@algocf
}
\makeatother

% ============================================================================
%                            TABEL (IEEE STYLE)
% ============================================================================
\usepackage{booktabs}        % Garis horizontal tabel berkualitas
\usepackage{siunitx}         % Perataan angka dan format numerik

% ============================================================================
%                          BIBLIOGRAFI (IEEE STYLE)
% ============================================================================
\usepackage{csquotes}
\usepackage[backend=biber,style=ieee,maxbibnames=99,minbibnames=99,sorting=none,isbn=false,url=false,date=year]{biblatex}
\addbibresource{references.bib}

\usepackage{microtype}
\usepackage{silence}

% ============================================================================
%                              HYPERREF
% ============================================================================
\usepackage[hidelinks]{hyperref}

% ============================================================================
%                          CLEVEREF (CROSS-REFERENCE)
% ============================================================================
% HARUS dimuat SETELAH hyperref
\usepackage[capitalize,noabbrev]{cleveref}
